Our second main finding is that LLMs have serious limitation on this task. The false positive rate on our audit of highest rated LLM validated transcripts was \data{summary_stats/tagging_performance/human_audit_false_positive_rate_pct_percentage}. Thus, while LLMs can be useful in identifying collusive communication, they can only produce high reliability results when used in conjunction with humans. For example, a plaintiff law firm using this methodology for discovery would likely need extensive human review to obtain useful results. At the same time, clearly the LLMs can greatly enhance the efficiency of this process, as evidenced by the large number of transcripts and hours of communication that we reviewed.

The key issue with the false positives is that collusive communications are rare. The overall percentage of validated transcripts is \data{summary_stats/tagging_performance/llm_validation_rate_overall_sample_pct_percentage}. Given the well-known problem that LLMs have occasional hallucinations, this means that many of the transcripts that are flagged end up being false positives.

Another interesting point is that the overall ability to detect collusive communication can be considerably increased by employing more human labor. In our study, out of all transcripts that were human-rated as collusive in the benchmark sample, only \data{summary_stats/tagging_performance/benchmark_collusive_validated_pct_percentage} were rediscovered in the LLM-validated sample, that only had \data{summary_stats/tagging_performance/llm_validation_tagged_int} transcripts. However, in a commercial application like legal discovery, once could choose to have humans go through the entire LLM-flagged sample. This would involve considerably more human labor, as that sample has \data{summary_stats/tagging_performance/llm_tagged_count_int} transcripts, and is likely to have a higher false positive rate. However, a full \data{summary_stats/tagging_performance/benchmark_collusive_flagged_pct_percentage} of the collusive communications in the benchmark sample was rediscovered in that sample. Thus, one could discover considerably more collusive communication at the cost of more human labor.

Finally, a caveat to our analysis is that we used a simple pipeline. Our large scale analysis makes a single LLM call to each transcript, followed by repetitions only in the top hits. Another avenue to increase performance are engineering improvements. It is likely that agentic systems, that make more calls, and decide what calls to make based on previous results, would be able to improve performance. While we have not explored this approach, it is currently being found useful in some of the most commercially successful LLM applications such as code generation.